\section{Clebsch-Gordan Tensor Product Details}
\label{sec:CGTP_details}
The most natural map is $V\times W\to V\otimes W$ constructed by taking an outer product of the inputs. If the inputs are explicitly written as a direct sum of irreps, we can write the tensor product as
\begin{align}
\label{eqn:tensor_product_full}
\x \otimes \y = \bigoplus_{\substack{\xl{1} \in \x \\ \yl{2} \in \y}}(\xl{1} \otimes \yl{2})
\end{align}
a new basis which is the sum of tensor product reps.

The key idea of a Clebsch-Gordan tensor product is we can explicitly reduce the tensor product reps back into a direct sum of irreps with a change of basis. This change of basis is the definition of the Clebsch-Gordan coefficients, giving us
\begin{align}
\label{eqn:cgtp_sum}
\xl{1}\otimes\yl{2} = & \bigoplus_{\ell_3} (\xl{1} \cgtp \yl{2})^{(\ell_3)}
\end{align}
where
\begin{align}
\label{eqn:cgtp}
&(\xl{1} \cgtp \yl{2})^{(\ell_3)}_{m_3} \nonumber \\
= &\sum_{m_1=-\ell_1}^{\ell_1}\sum_{m_2=-\ell_2}^{\ell_2} C^{\ell_3,m_3}_{\ell_1,m_1,\ell_2,m_2}\xl{1}_{m_1}\yl{2}_{m_2}.
\end{align}

Therefore the original tensor product can also be rewritten as a direct sum of irreps. This defines the Clebsch-Gordan tensor product (CGTP)
\begin{align}
\label{eqn:cgtp_full}
\x \cgtp \y = \bigoplus_{\substack{\xl{1} \in \x \\ \yl{2} \in \y}}(\xl{1} \cgtp \yl{2}).
\end{align}

Note that full CGTP is really just a change of basis from a sum of tensor product reps to a sum of irreps. Hence, we \textbf{do not lose any information}.

